% !TeX root = ../main.tex
\section{问题分析}

\subsection{问题一分析}

问题一的本质是在真实次品率未知的条件下，基于样本证据对“是否超过标称值”作出接收或拒收判断。由于两类信度方向相反，需分别构造方向相反的单侧检验，并把题目给出的信度转化为显著性水平；抽样机制决定检验分布，总体量未知或抽样比例较小时采用二项分布，已知有限总体且无放回抽样时采用超几何分布。通过枚举样本量与临界次品数，可在两类风险约束下确定最小可执行的固定样本方案；当样本证据不足以支持任一方向时，保留证据不足区域，并通过临界点相邻概率与参数敏感性检验方案的稳定性。

\subsection{问题二分析}

问题二属于带部分可观测质量信息和循环返工的序贯决策问题。零配件未检测时，其真实质量不可直接观察；成品一旦被判定为不合格，又会对两个零配件的质量提供联合信息。此时两个零配件的后验质量不再独立，若只记录边际次品率或预先固定一组检测变量，就会丢失“不合格成品”这一观测带来的相关性，也无法根据回收件的当前信息调整动作。

因此，本文以两类零配件质量的联合后验分布作为信念状态，把购件、检测零配件、装配、检测成品及拆解作为可行动作。每次检测或装配结果出现后，先按条件概率更新信念，再由 Bellman 方程比较当前动作成本与后续最小期望成本。求解输出不是一组贯穿全流程的固定 0--1 变量，而是状态到动作的映射。问题一给出的批次质量判定可用于确定问题二的次品率输入；若批次落在证据不足区，则应保留参数不确定性，而不能把标称次品率当作已确认事实。

\subsection{问题三分析}

问题三把问题二的单级装配扩展为两道工序、8 个零配件、3 个半成品和 1 个成品的树状结构。决策仍然是在每个环节选择是否检测、对不合格品是否拆解，但环节数量增至 16 个二元决策位，且拆解回流会跨越装配层级。若沿用问题二的信念状态方法，状态空间将随零部件数量急剧增长，难以直接枚举全部质量信息组合。

因此，问题三针对“给定一个固定策略”计算期望成本，转而建立有限状态马尔可夫奖励过程：以各零配件和半成品的真实质量状态构成系统状态，按质量传递规则和成本参数写出一步期望成本与状态转移，联立线性方程组求出交付一件合格成品的期望成本。固定策略的评价可以精确求解，且 $2^{16}=65536$ 个策略仍在可精确遍历范围内，故对全部策略进行枚举并直接获得固定策略空间内的全局最优解；遗传算法作为策略空间扩大后的扩展求解方法，Monte Carlo 模拟用于复核递推实现。问题三与问题二的区别在于：问题二强调信息状态相关策略，问题三则在固定策略空间中寻求可精确求解的全局最优方案，并以明确的策略类限制换取大规模结构的可计算性。

\subsection{问题四分析}

问题四的核心是把“次品率是抽样估计值”这一不确定性从问题一带入问题二与问题三。点估计口径只保留样本比例一个数值，忽略了小样本下真实次品率可能明显偏离样本比例的风险；问题一已经表明，同样的样本证据既可能支持接收、也可能支持拒收或证据不足，因此直接把表~1、表~2 的数值当作已知常数，会低估不检测策略的隐性成本。

处理思路是把每个次品率转化为后验分布：将名义值解释为 $n$ 件样本中的次品比例，在无信息均匀先验下得到 Beta 后验，再从后验抽取情景并按情景平均期望利润。这样既保留问题二、问题三的质量传递与成本规则（内层期望可解析计算），又使参数不确定性以期望、标准差与分位数等可量化指标进入决策；决策准则为风险中性的后验期望利润最大化，风险指标只作辅助描述。题面未给出实际抽样样本量，本问把 $n$ 作为证据强度参数并以 $n=40$ 为代表性小样本情景给出数值结果。由于状态相关信念在连续参数不确定下的完整求解规模很大，本问在问题二采用与 5.2.5 节一致的七变量固定策略类，在问题三沿用 16 位策略空间并以全枚举加遗传算法求解；完整信念状态 MDP 的参数不确定性版本作为模型局限在评价部分说明。

验证思路来自后验的收缩性质：样本量 $n$ 增大时，Beta 后验向名义值集中，后验最优策略与期望利润应收敛到问题二、问题三的点估计口径。若小样本下策略发生切换而大样本下恢复，则说明切换确由参数不确定性引起，而非数值噪声。

\subsection{各问之间的关系}

问题一解决采购批次能否接收以及质量信息可信度的问题，为生产阶段提供可用的次品率信息。问题二在给定或更新后的次品率基础上，把一次装配扩展为包含检测、市场损失、拆解和再次装配的信念状态决策；问题三进一步把装配结构扩展为多工序、多零配件，以期望递推和全策略精确枚举获得固定策略空间内的全局最优方案，并以遗传算法验证大规模扩展求解能力。前三问分别处理“参数是否可信”“给定质量参数后如何在小规模流程中行动”和“如何在大规模策略空间中优化”，共同形成从质量判断到生产处置、再到全流程优化的递进链条。问题四在前三问基础上把次品率还原为抽样估计值，用后验情景平均重新完成问题二与问题三的决策，使“参数是否可信”与“如何行动、如何优化”闭环：点估计口径给出基准方案，后验口径给出考虑抽样不确定性的修正方案，样本量增大时两者收敛。
